\chapter{Insights par Département}

\section{Analyse Comparative Inter-Départements}

L'attrition n'est pas un phénomène uniforme. Elle se concentre dans des poches spécifiques de l'organisation. Ce chapitre détaille le profil de risque pour chaque pôle d'activité.

\section{Le Département des Ventes (Sales) : Un Point d'Alerte Majeur}

Le pôle Sales affiche les indicateurs les plus préoccupants de toute l'entreprise.

\begin{table}[h]
\centering
\begin{tabular}{|l|c|c|}
\hline
\textbf{Indicateur} & \textbf{Sales} & \textbf{Moyenne Entreprise} \\ \hline
Taux d'Attrition & \textbf{31\%} & 23\% \\ \hline
Heures Moyennes & 225h & 201h \\ \hline
Satisfaction & 0.48 & 0.61 \\ \hline
\end{tabular}
\caption{Comparaison des indicateurs clés pour le département Sales}
\end{table}

\textbf{Analyse des causes :}
\begin{itemize}
    \item \textbf{Pression sur les objectifs :} La charge horaire est 12\% plus élevée que la moyenne.
    \item \textbf{Inadéquation Salariale :} Malgré une charge de travail supérieure, le salaire moyen en Sales est 15\% plus bas que dans les pôles techniques, créant un sentiment de déclassement.
\end{itemize}

\section{L'Ingénierie (Engineering) : Un Pilier de Stabilité}

À l'opposé des Sales, le département Engineering parvient à retenir ses talents avec un taux d'attrition de seulement \textbf{14\%}.

\textbf{Les facteurs de succès :}
\begin{itemize}
    \item \textbf{Équilibre Vie-Travail :} La charge mensuelle est stabilisée autour de 190h.
    \item \textbf{Reconnaissance :} Le taux de promotion dans les 5 dernières années est de 12\% (contre 4\% en Sales).
\end{itemize}

\section{Le Département IT : Une Rétention Basée sur la Qualité des Projets}

L'IT présente une situation hybride : une satisfaction élevée (\textbf{0.68}) mais une charge de travail parfois fluctuante. L'analyse montre que les départs en IT sont moins liés au stress qu'au manque de diversité dans les projets technologiques assignés (\textit{number\_project}).

\section{Le Département RH : "Le Cordonnier Mal Chaussé"}

Une observation ironique ressort de l'analyse : le département RH lui-même subit une attrition de \textbf{21\%}.
\begin{itemize}
    \item \textbf{Cause Identifiée :} Une surcharge liée justement à la gestion administrative du turnover des autres départements.
    \item \textbf{Impact :} Cette situation fragilise la capacité de l'entreprise à mener à bien ses politiques de rétention, créant un cercle vicieux.
\end{itemize}

\section{Matrice de Risque par Département}

Nous avons synthétisé ces insights dans une matrice de risque stratégique :

\begin{table}[h]
\centering
\begin{tabular}{|l|l|l|}
\hline
\textbf{Niveau de Risque} & \textbf{Départements} & \textbf{Priorité d'Action} \\ \hline
\cellcolor[HTML]{FFCCCB}Critique & Sales, Marketing & Immédiate (Mois 1) \\ \hline
\cellcolor[HTML]{FFFFE0}Modéré & HR, Operations & Trimestrielle (Mois 3) \\ \hline
\cellcolor[HTML]{90EE90}Faible & Engineering, IT & Surveillance (Mois 6) \\ \hline
\end{tabular}
\caption{Classification stratégique des départements par niveau de risque}
\end{table}
